% Математическая модель мультикоптера в Gazebo Sim (gz-sim).
% Сборка: pdflatex gazebo_drone_physics_model.ru.tex
%
% Источники:
%   https://github.com/gazebosim/gz-sim
%   MulticopterMotorModel.cc, LiftDrag.cc, Physics.cc
%   DART (движок по умолчанию), SDFormat 1.11

\documentclass[a4paper,11pt]{article}

\usepackage[T2A]{fontenc}
\usepackage[utf8]{inputenc}
\usepackage[russian]{babel}
\usepackage{amsmath,amssymb,amsfonts}
\usepackage{geometry}
\usepackage{hyperref}
\usepackage{enumitem}
\usepackage{array}
\usepackage{bm}
\geometry{left=2.5cm,right=2.5cm,top=2.5cm,bottom=2.5cm}
\setlist{nosep,leftmargin=*}
\hypersetup{colorlinks=true,linkcolor=blue,urlcolor=blue,citecolor=blue}

\newcommand{\code}[1]{\texttt{\detokenize{#1}}}
\newcommand{\R}{\mathbb{R}}
\newcommand{\vect}[1]{\bm{#1}}
\newcommand{\mat}[1]{\bm{#1}}
\newcommand{\norm}[1]{\left\lVert #1 \right\rVert}
\newcommand{\abs}[1]{\left\lvert #1 \right\rvert}

\title{Математическая модель мультикоптера\\
в симуляторе Gazebo Sim (\code{gz-sim})}
\author{}
\date{}

\begin{document}
\maketitle

\begin{abstract}
Документ описывает физическую и математическую модель квадрокоптера
в \textbf{Gazebo Sim} (\code{gz-sim}) и сопоставляет её с моделью
\textbf{ArduPilot SITL}, используемой в проекте \code{Drone\_Swarm\_Simulator\_v2}.
Даны уравнения движка \textbf{DART} (обобщённые координаты, матрица масс,
алгоритм ABA, LCP для контактов), полная цепочка обнаружения столкновений
(broadphase, GJK, EPA, mesh/convex hull), формулы контактных сил и трения,
модель ротора Gazebo и модель \code{SIM\_Motor} в SITL.
Показано, что в Gazebo для тех же кинематических величин
($\vect{p},\mat{R},\vect{v},\vect{\omega}$) решается существенно более
тяжёлая система: многосвязное тело, неявная интеграция, сотни контактных
ограничений при mesh-коллизиях, тогда как SITL в режиме \code{SITL-only}
интегрирует одно жёсткое тело без столкновений с окружением.
\end{abstract}

\tableofcontents
\newpage

%====================================================================
\section{Архитектура модели в Gazebo Sim}
%====================================================================

Gazebo Sim --- это симулятор робототехники, в котором динамика вычисляется
\textbf{физическим движком} (по умолчанию \textbf{DART} из пакета
\code{gz-physics}), а специализированные эффекты (двигатели, подъёмная сила
крыльев, ветер) добавляются \textbf{системными плагинами} (\emph{systems}).

Для типичного квадрокоптера (пример \code{examples/worlds/quadcopter.sdf})
цепочка выглядит так:
\begin{enumerate}
  \item \textbf{SDF-модель} задаёт геометрию, массу, инерцию, сочленения
  (joints), коллизионные формы и параметры поверхности.
  \item \textbf{Physics system} (\code{gz-sim-physics-system}) создаёт
  в DART скелет (\emph{skeleton}) из звеньев (\emph{links}) и суставов,
  включает гравитацию, интегрирует уравнения движения.
  \item На каждый ротор вешается \textbf{MulticopterMotorModel} --- по
  команде угловой скорости вычисляет тягу, реактивный момент и
  поперечное сопротивление винта; прикладывает силы/моменты к звеньям.
  \item Опционально: \textbf{MulticopterVelocityControl} (высокоуровневый
  закон скорости), \textbf{LiftDrag} (аэродинамика корпуса/крыльев),
  \textbf{Wind} (поле ветра).
  \item При пересечении коллизионных оболочек DART формирует
  \textbf{контактные ограничения}; решатель LCP выдаёт импульсы,
  предотвращающие проникновение и учитывающие трение.
\end{enumerate}

\paragraph{Связь с проектом Kursov3.}
В текущем репозитории \code{Drone\_Swarm\_Simulator\_v2} для полётной
динамики используется связка \textbf{ArduPilot SITL + Webots}, а не Gazebo.
Тем не менее Gazebo Sim --- стандартная среда для SITL с плагином
\code{ardupilot\_gazebo}; математическая структура модели (жёсткое тело +
квадратичная тяга роторов) совпадает с классической схемой PX4/Rotors/Gazebo.

%====================================================================
\section{Что такое SDF}
\label{sec:sdf}
%====================================================================

\textbf{SDF} (Simulation Description Format, \url{https://sdformat.org}) ---
это XML-язык описания робота и мира для Gazebo.
Файлы \code{*.sdf} --- не программа, а \textbf{паспорт модели}: откуда
физический движок берёт массы, формы, суставы и параметры контакта.

\subsection{Иерархия файла}

Типичная структура:
\begin{itemize}
  \item \code{world} --- гравитация, освещение, физика (\code{physics}),
  статические объекты (пол, стены);
  \item \code{model} --- робот (дрон), состоит из \code{link} и \code{joint};
  \item \code{link} --- одно \textbf{звено} (корпус, луч ротора, пропеллер);
  \item \code{joint} --- \textbf{сустав} между двумя \code{link};
  \item внутри \code{link}: \code{inertial}, \code{visual}, \code{collision}.
\end{itemize}

Gazebo при загрузке SDF строит дерево DART: каждый \code{link}
$\to$ \code{BodyNode}, каждый \code{joint} $\to$ \code{Joint}.

\subsection{Чем SDF отличается от URDF}

URDF --- формат ROS; SDF --- родной для Gazebo, поддерживает мир,
несколько моделей, мягкий контакт, плагины (\code{plugin}), ветер.
Для дрона в Gazebo обычно один \code{model} с 5--9 \code{link}
и 4--8 \code{joint}.

\subsection{Что задаётся числами в SDF}

\begin{center}
\renewcommand{\arraystretch}{1.2}
\begin{tabular}{|l|l|}
\hline
\textbf{Блок SDF} & \textbf{Влияние на уравнения} \\
\hline
\code{inertial/mass}, \code{inertia} & $m$, $\mat{I}$ в \eqref{eq:newton} \\
\code{collision/geometry} & форма для \eqref{eq:broadphase}--\eqref{eq:penetration} \\
\code{surface/friction} & $\mu$ в \eqref{eq:coulomb} \\
\code{surface/bounce} & $e$ в \eqref{eq:restitution} \\
\code{joint/type} & число DOF в $\vect{q}$ \\
\code{plugin MulticopterMotorModel} & $k_f$, $k_m$, $\tau_\uparrow$ \\
\hline
\end{tabular}
\end{center}

Без SDF симулятор не знает, \emph{какой} дрон считать: те же уравнения DART,
но другие параметры и топология.

%====================================================================
\section{Системы координат и состояние}
%====================================================================

\subsection{Мировая система координат (World, ENU)}

Gazebo по умолчанию использует правую декартову систему \textbf{ENU}
(East--North--Up):
\begin{itemize}
  \item ось $X$ --- на восток, м;
  \item ось $Y$ --- на север, м;
  \item ось $Z$ --- вверх, м.
\end{itemize}

Положение центра масс (ЦМ) корпуса:
\[
  \vect{p} = (p_x, p_y, p_z)^\top \in \R^3.
\]
Ориентация задаётся кватернионом $\vect{q} = (q_w, q_x, q_y, q_z)^\top$
или матрицей поворота $\mat{R}_{WB} \in SO(3)$ из связанной (body) системы
в мировую.

\subsection{Связанная система (Body)}

Оси корпуса жёстко закреплены на раме дрона. Вектор силы тяги $i$-го
ротора в локальной системе винта обычно направлен вдоль локальной оси $Z$
(вверх относительно диска пропеллера). В мировую систему:
\[
  \vect{F}_{T,i}^W = \mat{R}_{WB}\,\mat{R}_{B i}\,
  \begin{pmatrix} 0 \\ 0 \\ T_i \end{pmatrix},
\]
где $T_i$ --- скалярная тяга (Н), $\mat{R}_{B i}$ --- ориентация
звена ротора относительно корпуса.

\subsection{Полное состояние жёсткого тела (6-DOF)}

Минимальный набор для интегратора DART:
\begin{align*}
  \text{положение ЦМ:} &\quad \vect{p}(t), \\
  \text{ориентация:} &\quad \mat{R}_{WB}(t)\ \text{или}\ \vect{q}(t), \\
  \text{линейная скорость (мир):} &\quad \vect{v}(t) = \dot{\vect{p}}, \\
  \text{угловая скорость (связанная):} &\quad \vect{\omega}^B(t) \in \R^3.
\end{align*}
Эквивалентная форма Коши --- \textbf{12 скалярных уравнений первого порядка}
(3+3+3+3).

%====================================================================
\section{Уравнения движения жёсткого тела}
%====================================================================

DART решает \textbf{уравнения Ньютона--Эйлера} для каждого звена,
связанного с другими суставами или контактами.

\subsection{Линейное движение}

Сумма всех сил в мировой системе $\vect{F}^W$ даёт
\begin{equation}
  m\,\dot{\vect{v}} = \vect{F}^W
  = \vect{F}_g + \sum_i \vect{F}_{T,i}^W + \sum_j \vect{F}_{\mathrm{drag},j}^W
  + \vect{F}_{\mathrm{contact}} + \vect{F}_{\mathrm{ext}},
  \label{eq:newton}
\end{equation}
где:
\begin{itemize}
  \item $m$ --- масса звена (кг), из SDF \code{inertial/mass};
  \item $\vect{F}_g = m\vect{g}$ --- сила тяжести (Н);
  \item $\vect{F}_{T,i}^W$ --- тяга $i$-го ротора (Н);
  \item $\vect{F}_{\mathrm{drag},j}^W$ --- аэродинамическое сопротивление
  (модель винта или \code{LiftDrag});
  \item $\vect{F}_{\mathrm{contact}}$ --- контактная сила от DART (Н);
  \item $\vect{F}_{\mathrm{ext}}$ --- прочие внешние силы (плагины).
\end{itemize}

\subsection{Вращательное движение}

В связанной системе, относительно центра масс:
\begin{equation}
  \mat{I}^B \dot{\vect{\omega}}^B
  + \vect{\omega}^B \times (\mat{I}^B \vect{\omega}^B)
  = \vect{\tau}^B,
  \label{eq:euler}
\end{equation}
где:
\begin{itemize}
  \item $\mat{I}^B \in \R^{3\times 3}$ --- тензор инерции в связанных осях
  (кг\,м$^2$);
  \item $\vect{\omega}^B = (\omega_x,\omega_y,\omega_z)^\top$ --- угловая
  скорость (рад/с);
  \item $\vect{\tau}^B$ --- суммарный момент сил относительно ЦМ (Н\,м).
\end{itemize}

Момент складывается из:
\begin{equation}
  \vect{\tau}^B =
  \sum_i \left(\vect{r}_i^B \times \vect{F}_{T,i}^B\right)
  + \sum_i \vect{\tau}_{Q,i}^B
  + \vect{\tau}_{\mathrm{drag}}^B
  + \vect{\tau}_{\mathrm{contact}}^B,
  \label{eq:torque_sum}
\end{equation}
где $\vect{r}_i^B$ --- плечо $i$-го ротора от ЦМ, $\vect{\tau}_{Q,i}^B$ ---
реактивный (drag) момент пропеллера.

\subsection{От ускорения к новому состоянию}

На шаге $\Delta t$ DART использует \textbf{неявную (implicit) схему}
с учётом ограничений (контакты, суставы). В упрощённом виде (без контактов):
\begin{align}
  \vect{v}(t+\Delta t) &= \vect{v}(t) + \dot{\vect{v}}\,\Delta t,
  \label{eq:vel_int}\\
  \vect{p}(t+\Delta t) &= \vect{p}(t) + \vect{v}(t+\Delta t)\,\Delta t,
  \label{eq:pos_int}\\
  \dot{\mat{R}}_{WB} &= \mat{R}_{WB}\,[\vect{\omega}^B]_\times,
  \label{eq:rot_kin}
\end{align}
где $[\vect{a}]_\times$ --- кососимметрическая матрица:
\[
  [\vect{a}]_\times =
  \begin{pmatrix}
    0 & -a_z & a_y \\
    a_z & 0 & -a_x \\
    -a_y & a_x & 0
  \end{pmatrix}.
\]

\textbf{Цепочка вывода:}
команда ротора $\Omega_i^{\mathrm{ref}}$
$\Rightarrow$ фильтр $\Rightarrow$ $\omega_i$
$\Rightarrow$ тяга $T_i$
$\Rightarrow$ $\vect{F}, \vect{\tau}$
$\Rightarrow$ \eqref{eq:newton}, \eqref{eq:euler}
$\Rightarrow$ $\dot{\vect{v}}, \dot{\vect{\omega}}$
$\Rightarrow$ \eqref{eq:vel_int}--\eqref{eq:rot_kin}
$\Rightarrow$ новое $(\vect{p}, \mat{R}, \vect{v}, \vect{\omega})$.

%====================================================================
\section{Движок DART: обобщённые координаты и динамика}
%====================================================================

Gazebo Sim по умолчанию передаёт динамику библиотеке \textbf{DART}
(Dynamic Animation and Robotics Toolkit, \url{https://github.com/dartsim/dart}).
В отличие от SITL, где состояние --- одно жёсткое тело, DART описывает
\textbf{скелет} (\emph{skeleton}) --- дерево звеньев (\code{BodyNode}),
соединённых суставами (\code{Joint}) с обобщёнными координатами $\vect{q}$.

\subsection{Обобщённые координаты и степени свободы}

Для квадрокоптера в SDF типична структура:
\begin{itemize}
  \item корпус (\code{base\_link}) --- 6 DOF относительно мира
  (если \code{free joint}) или 0 DOF (если закреплён);
  \item 4 ротора --- по 1 DOF (\code{revolute joint}) на каждый винт.
\end{itemize}
Вектор обобщённых координат:
\[
  \vect{q} =
  \bigl(
    \vect{p}^\top,\ \vect{\eta}^\top,\ \theta_0,\ldots,\theta_{n-1}
  \bigr)^\top,
\]
где $\vect{\eta}$ --- параметры ориентации корпуса (кватернион или углы),
$\theta_i$ --- углы поворота суставов роторов.
Скорости $\dot{\vect{q}} = \vect{u}$; ускорения $\ddot{\vect{q}} = \vect{a}$.

\subsection{Уравнения динамики в обобщённых координатах}

DART записывает динамику скелета в форме Лагранжа:
\begin{equation}
  \mat{M}(\vect{q})\,\ddot{\vect{q}}
  + \vect{C}(\vect{q},\dot{\vect{q}})
  + \vect{G}(\vect{q})
  = \vect{\tau} + \vect{\tau}_{\mathrm{ext}} + \vect{J}_c^\top \vect{\lambda},
  \label{eq:dart_eom}
\end{equation}
где:
\begin{itemize}
  \item $\mat{M}(\vect{q}) \in \R^{n\times n}$ --- \textbf{матрица обобщённых масс}
  (кг\,м$^2$ или эквивалентные единицы по DOF);
  \item $\vect{C}$ --- вектор кориолисовых и центробежных сил;
  \item $\vect{G}$ --- обобщённые силы тяжести;
  \item $\vect{\tau}$ --- управляющие моменты/силы суставов (от моторов);
  \item $\vect{\tau}_{\mathrm{ext}}$ --- внешние силы (плагины Gazebo);
  \item $\vect{J}_c$ --- якобиан контактных/кинематических ограничений;
  \item $\vect{\lambda}$ --- множители Лагранжа (импульсы контактов).
\end{itemize}

\paragraph{Связь с Ньютоном--Эйлером.}
Для одного звена \eqref{eq:dart_eom} эквивалентна
\eqref{eq:newton}--\eqref{eq:euler}, но $\mat{M}$, $\vect{C}$, $\vect{G}$
собираются с учётом кинематической цепи (плеч роторов, вращающихся звеньев).

\subsection{Алгоритм ABA (Featherstone)}

Без контактов DART вычисляет ускорения за $O(n)$ операций
(\emph{Articulated Body Algorithm}):
\begin{equation}
  \ddot{\vect{q}} =
  \mat{M}(\vect{q})^{-1}
  \bigl(
    \vect{\tau} + \vect{\tau}_{\mathrm{ext}}
    - \vect{C}(\vect{q},\dot{\vect{q}})
    - \vect{G}(\vect{q})
  \bigr).
  \label{eq:aba}
\end{equation}
Два прохода по дереву звеньев:
\begin{enumerate}
  \item \textbf{от листьев к корню} --- накопление инерции сочленённого тела;
  \item \textbf{от корня к листьям} --- распространение ускорений.
\end{enumerate}

\subsection{Матрица масс и обратная динамика}

Матрица $\mat{M}$ может строиться алгоритмом CRBA
(\emph{Composite Rigid Body Algorithm}, сложность $O(n^2)$).
Обратная динамика (RNEA) по заданному $\ddot{\vect{q}}$ даёт необходимый $\vect{\tau}$.

\subsection{Пассивные силы суставов}

Каждый \code{Joint} в DART может иметь:
\begin{equation}
  \tau_j^{\mathrm{passive}} =
  -k_j (q_j - q_{j,0})
  - d_j \dot{q}_j
  - f_j \,\sgn(\dot{q}_j),
  \label{eq:joint_passive}
\end{equation}
где $k_j$ --- жёсткость, $d_j$ --- демпфирование, $f_j$ --- сухое трение.
В SDF задаётся через \code{physics/dart} или параметры сустава.

\subsection{Полный шаг \code{World::step()} в DART}

По документации архитектуры DART (\code{docs/onboarding/architecture.md}):
\begin{enumerate}
  \item \textbf{Прямая кинематика:}
  $\vect{q} \Rightarrow$ положения/скорости всех \code{BodyNode};
  \item \textbf{Прямая динамика (ABA):} $\ddot{\vect{q}}_{\mathrm{free}}$
  из \eqref{eq:aba} без контактов;
  \item \textbf{ConstraintSolver::solve():}
  \begin{enumerate}
    \item collision detection $\Rightarrow$ контакты;
    \item построение \code{ContactConstraint}, лимитов суставов;
    \item сбор boxed LCP; решение (Dantzig / PGS);
    \item применение импульсов:
    $\vect{v}_{n+1} = \vect{v}_n + \mat{M}^{-1}\vect{J}_c^\top \vect{\lambda}$.
  \end{enumerate}
  \item \textbf{Интеграция (semi-implicit Euler):}
  \begin{align}
    \dot{\vect{q}}_{n+1} &= \dot{\vect{q}}_n + \ddot{\vect{q}}_n\,\Delta t,
    \label{eq:dart_vel}\\
    \vect{q}_{n+1} &= \vect{q}_n + \dot{\vect{q}}_{n+1}\,\Delta t.
    \label{eq:dart_pos}
  \end{align}
\end{enumerate}

\paragraph{Отличие от SITL.}
SITL использует \textbf{явный} Эйлер для скорости и положения
(см.\ раздел~\ref{sec:sitl}); DART --- \textbf{полунеявный} Эйлер
\eqref{eq:dart_vel}--\eqref{eq:dart_pos} \emph{после} решения LCP,
что стабильнее при жёстких контактах.

%====================================================================
\section{Масса, центр масс и инерция (SDF)}
%====================================================================

В файле SDF для каждого \code{link} задаётся блок \code{inertial}:

\begin{itemize}
  \item \code{mass} --- $m$ (кг);
  \item \code{pose} --- смещение и поворот ЦМ относительно звена;
  \item \code{inertia} --- компоненты симметричного тензора
  $I_{xx}, I_{xy}, I_{xz}, I_{yy}, I_{yz}, I_{zz}$.
\end{itemize}

Матрица инерции в системе, заданной в SDF:
\[
  \mat{I} =
  \begin{pmatrix}
    I_{xx} & I_{xy} & I_{xz} \\
    I_{xy} & I_{yy} & I_{yz} \\
    I_{xz} & I_{yz} & I_{zz}
  \end{pmatrix}.
\]

Если оси SDF не совпадают с главными осями инерции, DART переносит
$\mat{I}$ в мировую/связанную систему через теорему Штейнера:
\begin{equation}
  \mat{I}^B = \mat{R}\,\mat{I}_{\mathrm{SDF}}\,\mat{R}^\top,
  \qquad
  \mat{I}_{\mathrm{parallel}} = \mat{I}_C + m\bigl(\norm{\vect{r}}^2\mat{E}
  - \vect{r}\vect{r}^\top\bigr),
  \label{eq:steiner}
\end{equation}
где $\vect{r}$ --- вектор от старого центра к новому, $\mat{E}$ --- единичная
матрица.

\paragraph{Практика для дрона.}
Корпус моделируют одним или несколькими звеньями; роторы --- отдельными
звеньями с малыми массами, соединёнными \code{revolute joint} (ось вращения
пропеллера). Суммарная масса и инерция всего аппарата --- результат
композиции всех звеньев в дереве кинематики DART.

%====================================================================
\section{Сила тяжести}
%====================================================================

В SDF мира (\code{world/physics/gravity}) задаётся вектор
\[
  \vect{g} = (g_x, g_y, g_z)^\top,
\]
по умолчанию $\vect{g} = (0,\,0,\,-9.80665)^\top$ м/с$^2$ в ENU
(направлен \emph{вниз}, т.\,е.\ отрицательная $Z$).

Сила тяжести на звено:
\begin{equation}
  \vect{F}_g = m\,\vect{g}.
  \label{eq:gravity}
\end{equation}

Physics system читает компонент \code{Gravity} и передаёт $\vect{g}$ в
движок (\code{world.SetGravity(...)}). Для отдельного звена гравитацию
можно отключить (\code{gravity\_enabled=false}), но для дрона она включена.

%====================================================================
\section{Модель одного двигателя: MulticopterMotorModel}
%====================================================================

Системный плагин \code{gz::sim::systems::MulticopterMotorModel}
(исходник \code{src/systems/multicopter\_motor\_model/}) реализует
классическую модель из пакета Rotors/PX4 (Furrer et al., ETH; далее
Martin \& Salaün, ICRA 2010 для поперечного потока).

\subsection{Вход: команда угловой скорости}

На топик \code{commandSubTopic} (например \code{gazebo/command/motor\_speed})
приходит сообщение \code{Actuators} с массивом \code{velocity}.
Для ротора с индексом \code{actuator\_number} $= i$:
\begin{equation}
  \Omega_i^{\mathrm{ref}} =
  \min\bigl(\code{velocity}[i],\,\Omega_{\max}\bigr),
  \qquad
  \Omega_{\max} = \code{maxRotVelocity}\ \text{рад/с}.
  \label{eq:ref_input}
\end{equation}

\subsection{Фильтр первого порядка (инерция ESC+мотор)}

Реальная скорость ротора не мгновенна. Применяется дискретный фильтр
с разными постоянными времени на разгон и торможение
(\code{timeConstantUp}, \code{timeConstantDown}):
\begin{equation}
  \Omega_i[k+1] =
  \begin{cases}
    e^{-\Delta t/\tau_\uparrow}\,\Omega_i[k]
    + \bigl(1-e^{-\Delta t/\tau_\uparrow}\bigr)\,\Omega_i^{\mathrm{ref}},
    & \Omega_i^{\mathrm{ref}} > \Omega_i[k], \\[6pt]
    e^{-\Delta t/\tau_\downarrow}\,\Omega_i[k]
    + \bigl(1-e^{-\Delta t/\tau_\downarrow}\bigr)\,\Omega_i^{\mathrm{ref}},
    & \text{иначе}.
  \end{cases}
  \label{eq:motor_filter}
\end{equation}
Параметры:
\begin{itemize}
  \item $\Delta t$ --- шаг симуляции, с;
  \item $\tau_\uparrow = \code{timeConstantUp}$, $\tau_\downarrow =
  \code{timeConstantDown}$, с;
  \item $\Omega_i$ --- отфильтрованная угловая скорость, рад/с.
\end{itemize}

Непрерывный аналог: $ \dot{\Omega} = (\Omega^{\mathrm{ref}}-\Omega)/\tau$.

\subsection{Связь с joint и параметр slowdown}

Чтобы избежать численных проблем, \emph{скорость сустава} в DART
уменьшается в \code{rotorVelocitySlowdownSim} $= n_s$ раз (типично $n_s=10$):
\begin{equation}
  \omega_i^{\mathrm{joint}} = \frac{\sigma_i\,\Omega_i}{n_s},
  \label{eq:joint_vel}
\end{equation}
где $\sigma_i \in \{+1,-1\}$ --- \code{turningDirection}
(\code{cw} $\Rightarrow -1$, \code{ccw} $\Rightarrow +1$).

Для расчёта тяги используется \emph{полная} скорость:
\begin{equation}
  \omega_i^{\mathrm{real}} = \omega_i^{\mathrm{joint}}\, n_s
  = \sigma_i\,\Omega_i.
  \label{eq:real_vel}
\end{equation}
Т.\,е.\ $n_s$ сокращается: тяга зависит от $\Omega_i$, а не от
завышенной joint-скорости.

\subsection{Тяга одного ротора}

\begin{equation}
  T_i = \sigma_i\,\sgn(\omega_i^{\mathrm{real}})\,
  \bigl(\omega_i^{\mathrm{real}}\bigr)^2\, k_f,
  \qquad
  k_f = \code{motorConstant}\ \ \text{Н/(рад/с)}^2.
  \label{eq:thrust}
\end{equation}

Физический смысл: $T \propto \omega^2$ --- упрощённая модель
несущего винта (аналог $T \propto n^2$ для об/мин).
Сила прикладывается к звену ротора вдоль его локальной $Z$ в мировой
системе:
\[
  \vect{F}_{T,i}^W = \mat{R}_{\mathrm{rotor},i}\,
  \begin{pmatrix}0\\0\\T_i\end{pmatrix}.
\]

\subsection{Реактивный момент (drag torque)}

Момент сопротивления вращению пропеллера (от тяги):
\begin{equation}
  \vect{\tau}_{Q,i}^{\mathrm{rotor}} =
  \begin{pmatrix} 0 \\ 0 \\ -\sigma_i\, T_i\, k_m \end{pmatrix},
  \qquad
  k_m = \code{momentConstant}\ \ \text{м},
  \label{eq:drag_torque}
\end{equation}
затем переносится в систему родительского звена (корпуса) и в мир.

\subsection{Поперечное сопротивление винта (blade flapping proxy)}

По Martin \& Salaün (2010), с учётом ветра (\code{Wind}):
\begin{align}
  \vect{v}_{\mathrm{rel}}^W &= \vect{v}_{\mathrm{link}}^W - \vect{v}_{\mathrm{wind}}^W,
  \label{eq:vrel}\\
  \vect{v}_\perp &= \vect{v}_{\mathrm{rel}}^W
  - (\vect{v}_{\mathrm{rel}}^W \cdot \vect{e}_\omega)\,\vect{e}_\omega,
  \label{eq:vperp}\\
  \vect{F}_{\mathrm{blade},i}^W} &=
  -\abs{\omega_i^{\mathrm{real}}}\, b_f\,\vect{v}_\perp,
  \label{eq:blade_drag}
\end{align}
где $\vect{e}_\omega$ --- ось вращения ротора в мировой системе,
$b_f = \code{rotorDragCoefficient}$ (Н/(м/с)$^2$ в документации кода).

\subsection{Rolling moment}

Дополнительный момент от поперечного потока:
\begin{equation}
  \vect{\tau}_{\mathrm{roll},i}^W =
  -\abs{\omega_i^{\mathrm{real}}}\,\mu_1\,\vect{v}_\perp,
  \qquad
  \mu_1 = \code{rollingMomentCoefficient}.
  \label{eq:rolling}
\end{equation}

\subsection{Сводка по одному двигателю}

\begin{center}
\renewcommand{\arraystretch}{1.3}
\begin{tabular}{|l|l|l|}
\hline
\textbf{Параметр SDF} & \textbf{Символ} & \textbf{Единица} \\
\hline
\code{motorConstant} & $k_f$ & Н/(рад/с)$^2$ \\
\code{momentConstant} & $k_m$ & м \\
\code{maxRotVelocity} & $\Omega_{\max}$ & рад/с \\
\code{rotorDragCoefficient} & $b_f$ & см. код \\
\code{rollingMomentCoefficient} & $\mu_1$ & см. код \\
\code{timeConstantUp/Down} & $\tau_\uparrow, \tau_\downarrow$ & с \\
\code{rotorVelocitySlowdownSim} & $n_s$ & --- \\
\code{turningDirection} & $\sigma_i$ & $\pm 1$ \\
\hline
\end{tabular}
\end{center}

%====================================================================
\section{Полная модель квадрокоптера}
%====================================================================

\subsection{Геометрия X-конфигурации}

Для 4 роторов (пример X3 в \code{quadcopter.sdf}) индексы $i=0,1,2,3$.
Роторы 0 и 1 --- \code{ccw}, 2 и 3 --- \code{cw} (или наоборот),
чтобы суммарный реактивный момент мог компенсироваться.

Плечо ротора $\vect{r}_i^B = (r_{x,i}, r_{y,i}, 0)^\top$ --- расстояние
от ЦМ до оси $i$-го винта в плоскости рамы.

\subsection{Суммарная тяга и моменты}

\begin{align}
  T_{\Sigma} &= \sum_{i=0}^{3} T_i,
  \label{eq:total_thrust}\\
  \vect{F}_T^W &= \sum_{i=0}^{3} \vect{F}_{T,i}^W,
  \label{eq:total_force}\\
  \vect{\tau}^B &= \sum_{i=0}^{3}
  \vect{r}_i^B \times
  \begin{pmatrix}0\\0\\T_i\end{pmatrix}
  + \sum_{i=0}^{3} \vect{\tau}_{Q,i}^B
  + \sum_{i=0}^{3} \vect{\tau}_{\mathrm{roll},i}^B
  + \vect{\tau}_{\mathrm{ext}}^B.
  \label{eq:total_torque}
\end{align}

В \emph{идеализированной} модели (без $\vect{v}_\perp$, одинаковые $T_i$):
\begin{itemize}
  \item вертикальный подъём: $T_{\Sigma} \approx mg$ в висении;
  \item крен: разность тяг левых/правых роторов создаёт $\tau_x$;
  \item тангаж: разность передних/задних --- $\tau_y$;
  \item рыскание: разность $\sigma_i T_i k_m$ --- $\tau_z$
  (компensation yaw).
\end{itemize}

\subsection{MulticopterVelocityControl (опционально)}

Плагин \code{MulticopterVelocityControl} по желаемой линейной скорости
и yaw-rate в связанной системе решает обратную задачу и выдаёт
$\Omega_i^{\mathrm{ref}}$ в компонент \code{Actuators}. Сами по себе
силы он не прикладывает --- нужен \code{MulticopterMotorModel} на каждом
роторе (см.\ API Gazebo Sim 10).

%====================================================================
\section{Суставы (\code{joint}) дрона в Gazebo}
\label{sec:joints}
%====================================================================

В SDF сустав связывает \textbf{parent link} и \textbf{child link}.
DART превращает тип joint в ограничения на относительное движение
и добавляет соответствующие строки в $\vect{q}$.

\subsection{Типовая кинематическая цепь квадрокоптера}

Пример (модель X3 / iris-like):
\begin{center}
\begin{tabular}{c}
\code{world} \\
$\downarrow$ \code{free} или фикс. база \\
\code{base\_link} (корпус) \\
$\swarrow$ \code{revolute} $\times 4$ \\
\code{rotor\_0..3} (звенья пропеллеров)
\end{tabular}
\end{center}

\subsection{Типы суставов и число DOF}

\begin{center}
\renewcommand{\arraystretch}{1.3}
\begin{tabular}{|l|c|l|}
\hline
\textbf{SDF \code{joint type}} & \textbf{DOF} & \textbf{Роль на дроне} \\
\hline
\code{fixed} & 0 &
луч жёстко к корпусу (упрощённая модель без вращения винта) \\
\code{revolute} & 1 &
ось пропеллера: угол $\theta_i$, скорость $\dot{\theta}_i$
(\code{MulticopterMotorModel} задаёт $\dot{\theta}_i$ через
\eqref{eq:joint_vel}) \\
\code{continuous} & 1 &
как \code{revolute}, но без лимита $2\pi$ (полный оборот) \\
\code{ball} & 3 & шаровой шарнир (редко на коптере) \\
\code{free} & 6 &
весь дрон как одно звено относительно мира
($x,y,z$, roll, pitch, yaw) \\
\hline
\end{tabular}
\end{center}

\subsection{Уравнения для \code{revolute} (ротор)}

Ось $\vect{e}_i$ в системе parent. Угол $\theta_i$ --- координата в $\vect{q}$.
Кинематика:
\begin{equation}
  \mat{R}_{\mathrm{child}} =
  \mat{R}_{\mathrm{parent}}\,\mat{R}_z(\theta_i),
  \qquad
  \dot{\mat{R}}_{\mathrm{child}} =
  \mat{R}_{\mathrm{parent}}\,\mat{S}(\vect{e}_i)\,\dot{\theta}_i.
  \label{eq:revolute_kin}
\end{equation}
DART добавляет в $\mat{M}(\vect{q})$ инерцию звена ротора и связь
между $\ddot{\theta}_i$ и ускорением корпуса.

Управление от плагина --- не момент $\tau_i$, а \textbf{скорость}:
\begin{equation}
  \dot{\theta}_i^{\mathrm{cmd}} =
  \frac{\sigma_i\,\Omega_i}{n_s}
  \quad \text{(\code{JointVelocityCmd})}.
  \label{eq:joint_cmd}
\end{equation}
Тяга \eqref{eq:thrust} считается по $\Omega_i$, а не по реакции сустава.

\subsection{Сустав \code{free} (весь дрон)}

Если корпус соединён с миром через \code{free joint}:
\[
  \vect{q}_{\mathrm{body}} =
  (x,y,z,\phi,\theta,\psi)^\top \ \text{или кватернион},
\]
6 DOF. Тогда \eqref{eq:dart_eom} описывает полёт целиком;
роторы дают +4 DOF $\Rightarrow$ типично $n_q = 10$.

Если корпус \textbf{приклеен} к миру (\code{fixed} к world) --- 0 DOF
корпуса, только роторы крутятся (редко для полёта).

\subsection{Пассивные и ограничивающие силы сустава}

Из \eqref{eq:joint_passive} и лимитов:
\begin{equation}
  \tau_i^{\mathrm{limit}} =
  \begin{cases}
    \lambda_i \geq 0, & \theta_i = \theta_{\max}, \\
    \lambda_i \leq 0, & \theta_i = \theta_{\min}.
  \end{cases}
\end{equation}
Для \code{continuous} лимитов нет; для \code{revolute} с \code{limit}
--- удар об упор при превышении угла.

\subsection{Отличие от SITL}

В SITL нет отдельных $\theta_i$ в состоянии: четыре мотора --- это
алгебраические $T_i$, а не DOF дерева. Вращение пропеллера
не интегрируется как $\dot{\theta}_i$; только $\vect{\omega}$ всего аппарата
в \eqref{eq:sitl_omega}.

%====================================================================
\section{Дополнительная аэродинамика: LiftDrag}
%====================================================================

Для несущих поверхностей (крыло, фюзеляж) плагин \code{LiftDrag} добавляет
силы в точке \code{cp} (center of pressure):
\begin{align}
  q &= \tfrac{1}{2}\rho \norm{\vect{v}_{\mathrm{plane}}}^2,
  \quad \text{динамическое давление}, \label{eq:q}\\
  \alpha &= \text{угол атаки в плоскости profile}, \\
  C_L &= \code{cla}\cdot\alpha \ \text{(до stall)}, \\
  C_D &= \abs{\code{cda}\cdot\alpha}, \\
  \vect{F}_{\mathrm{lift}} &= C_L\, q\, S\,\vect{e}_{\mathrm{lift}}, \\
  \vect{F}_{\mathrm{drag}} &= C_D\, q\, S\,(-\vect{e}_{\mathrm{drag}}),
\end{align}
где $\rho = \code{air\_density}$, $S = \code{area}$.
Полная сила $\vect{F} = \vect{F}_{\mathrm{lift}}+\vect{F}_{\mathrm{drag}}$,
момент $\vect{\tau} = \vect{\tau}_C + \vect{r}_{cp}\times\vect{F}$.

Для классического квадрокоптера без крыльев этот плагин часто не
используют; для гибридных VTOL --- обязателен.

%====================================================================
\section{Внутренние и внешние силы: классификация}
%====================================================================

\begin{center}
\renewcommand{\arraystretch}{1.2}
\begin{tabular}{|p{3.2cm}|p{4.5cm}|p{6.5cm}|}
\hline
\textbf{Тип} & \textbf{Источник в Gazebo} & \textbf{Физическая природа} \\
\hline
Внешняя &
$\vect{F}_g$, ветер, контакт с землёй &
Гравитация, среда, ограничение среды \\
\hline
Внутренняя (управление) &
\code{MulticopterMotorModel} &
Тяга и момент от пропеллеров \\
\hline
Внутренняя (диссипация) &
\code{rotorDragCoefficient}, joint damping &
Сопротивление потоку, трение в осях \\
\hline
Связь / constraint &
DART joints, contacts &
Идеальные кинематические связи, непроникновение \\
\hline
\end{tabular}
\end{center}

\paragraph{«Внутренние» силы} --- генерируются actuators модели (моторы).
\paragraph{«Внешние»} --- действуют от окружения (гравитация, ветер, пол, стены).

%====================================================================
\section{Столкновения: геометрия, обнаружение, контактные силы}
\label{sec:collision}
%====================================================================

В \textbf{ArduPilot SITL-only} (режим проекта Kursov3) столкновений
с землёй и препятствиями \emph{нет}: высота не ограничена снизу физикой.
В \textbf{Gazebo + DART} контакт --- обязательная часть каждого шага
\code{World::step()}.

\subsection{Что такое \code{link} в коллизии}

\textbf{Link} (звено) --- это жёсткая часть модели с собственной системой
координат, массой и (опционально) коллизией.
В SDF:
\begin{verbatim}
<link name="base_link">
  <inertial>...</inertial>
  <visual>...</visual>      <!-- только картинка -->
  <collision name="col_base"> <!-- только физика -->
    <geometry><box>...</box></geometry>
    <surface>...</surface>
  </collision>
</link>
\end{verbatim}

Важно:
\begin{itemize}
  \item \textbf{Коллизия привязана к link}, не к «дрону целиком».
  У квадрокоптера могут быть отдельные \code{collision} на корпусе,
  на лучах, на защитных ободах; у каждого --- свой \code{CollisionObject} в DART.
  \item \code{visual} и \code{collision} \textbf{независимы}: можно рисовать
  детальную сетку, а для физики использовать упрощённый box.
  \item DART хранит форму на \code{ShapeNode}, висящем на \code{BodyNode}
  данного link. При контакте силы прикладываются к \textbf{тому link},
  чья форма пересеклась (через якобиан в \eqref{eq:lcp_A}).
\end{itemize}

Если дрон касается пола лучом и корпусом одновременно, в LCP появляются
\textbf{две группы} контактов с разными $\vect{c}_k$ и разными звеньями
в дереве --- отсюда «сложность mesh + несколько link».

\subsection{Какой тип столкновения реально используется}

В таблице ниже --- все типы \code{collision/geometry} в SDF.
На практике для мультикоптера:

\begin{center}
\renewcommand{\arraystretch}{1.25}
\begin{tabular}{|l|l|l|l|}
\hline
\textbf{Тип} & \textbf{Математика} & \textbf{Сложность} & \textbf{Где применяют} \\
\hline
\code{plane} &
$\mathcal{P}=\{\vect{x}: \vect{n}^\top\vect{x}+d=0,\ \vect{n}^\top\vect{x}\geq 0\}$
& низкая & пол мира, стены \\
\code{box} &
$\mathcal{B}=\{\vect{x}: |R^\top(\vect{x}-\vect{c})| \leq \vect{a}\}$
& низкая & упрощённая рама, стол \\
\code{sphere} &
$\|\vect{x}-\vect{c}\|\leq r$ & низкая & сфера-прокси \\
\code{cylinder} &
радиус $r$, ось $\vect{e}$ & низкая & мачта, столб \\
\code{mesh} (как есть) &
$\bigcup_{\triangle} \mathrm{conv}\{v_{i_0},v_{i_1},v_{i_2}\}$
& \textbf{высокая} & редко (нестабильно) \\
\code{mesh}+\code{convex\_hull} &
$\mathrm{conv}(\mathcal{V})$ & средняя & корпус iris \\
\code{mesh}+\code{convex\_decomposition} &
$\bigcup_k \mathrm{conv}(\mathcal{V}_k)$ & высокая & сложная рама \\
\hline
\end{tabular}
\end{center}

\paragraph{Почему не «сырой» mesh.}
Невыпуклая сетка (вогнутые карманы под мотором) даёт множественные
ложные контакты и нестабильный EPA. Поэтому Gazebo при загрузке часто
\textbf{заменяет} mesh на выпуклую оболочку или набор выпуклых частей
(V-HACD). Визуально дрон детальный, физически --- упрощённый полиэдр.

\paragraph{Типичный выбор для посадки.}
\begin{itemize}
  \item \textbf{Земля:} \code{plane} --- аналитический контакт, 1--мало точек;
  \item \textbf{Дрон:} \code{box} или \code{convex hull} корпуса;
  \item \textbf{Пропеллеры:} часто \code{collision enabled=false}, чтобы
  лопасти не цепляли воздух/землю ложными контактами.
\end{itemize}

\paragraph{Сетка (mesh) и выпуклая оболочка.}
Исходный \code{.dae}/\code{.stl} содержит вершины
$\mathcal{V}=\{\vect{v}_0,\ldots,\vect{v}_{N-1}\}$ и грани
$\mathcal{F}=\{f_0,\ldots,f_{M-1}\}$.
Для устойчивой физики Gazebo рекомендует не использовать произвольную
невыпуклую сетку напрямую, а:
\begin{enumerate}
  \item \textbf{Convex hull} (\code{convex\_hull}): минимальная выпуклая
  оболочка $\mathrm{conv}(\mathcal{V})$;
  \item \textbf{Convex decomposition} (V-HACD, \code{convex\_decomposition}):
  разбиение на $K$ выпуклых субмешей
  $\mathcal{M}=\bigcup_{k=1}^{K} \mathrm{conv}(\mathcal{V}_k)$.
\end{enumerate}
Каждый субмеш --- отдельный \code{CollisionObject} в DART/FCL.

\subsection{Двухфазное обнаружение столкновений}

\subsubsection{Широкая фаза (broadphase) --- подробно}

\textbf{Задача:} из $m$ коллизионных объектов (все \code{collision} всех
\code{link} всех моделей + пол) найти пары, которые \emph{могут}
пересекаться, не проверяя каждую пару точно ($O(m^2)$ дорого).

\paragraph{Ограничивающий ящик AABB.}
Для объекта $i$ в мировой системе:
\begin{equation}
  \mathrm{AABB}_i =
  \bigl\{
    \vect{x} :
    x_{\min,i} \leq x \leq x_{\max,i},\;
    y_{\min,i} \leq y \leq y_{\max,i},\;
    z_{\min,i} \leq z \leq z_{\max,i}
  \bigr\},
  \label{eq:aabb}
\end{equation}
где границы берутся по всем точкам формы link в текущей позе
(при повороте дрона AABB \textbf{расширяется}, становится больше реального тела).

\paragraph{Тест пересечения AABB.}
Два ящика пересекаются тогда и только тогда, когда по \emph{каждой} оси
интервалы проекций перекрываются:
\begin{equation}
  \mathrm{AABB}_i \cap \mathrm{AABB}_j \neq \emptyset
  \Leftrightarrow
  \forall \alpha \in \{x,y,z\}:\;
  \max(a_{\min,i}^{(\alpha)}, a_{\min,j}^{(\alpha)})
  \leq
  \min(a_{\max,i}^{(\alpha)}, a_{\max,j}^{(\alpha)}).
  \label{eq:aabb_test}
\end{equation}

\paragraph{Множество кандидатов.}
\begin{equation}
  \mathcal{P}_{\mathrm{cand}} =
  \bigl\{
    (i,j) : i<j,\;
    \mathrm{AABB}_i \cap \mathrm{AABB}_j \neq \emptyset,\;
    \text{слои } i,j \text{ не исключены фильтром}
  \bigr\}.
  \label{eq:broadphase}
\end{equation}
Фильтр \code{collide\_bitmask} в SDF отключает, например, контакт
«корпус--лопасть того же дрона».

\paragraph{Ускорение: BVH / sweep-and-prune.}
На каждом шаге $\Delta t$ позы меняются мало; используют дерево
ограничивающих объёмов (BVH) или сортировку по оси
(\emph{sweep-and-prune}): сложность порядка $O(m \log m)$--$O(m)$,
а не $O(m^2)$.

\paragraph{Связь с link.}
Один \code{link} может иметь \textbf{несколько} \code{collision}
(например, корпус + обод мотора) --- это \textbf{отдельные} объекты
$i,j$ в broadphase. Поэтому число $m$ в формулах больше числа дронов.

\subsubsection{Узкая фаза (narrowphase)}

Для каждой пары $(i,j) \in \mathcal{P}_{\mathrm{cand}}$ вычисляется
пересечение точных форм. В Gazebo (через FCL) для выпуклых тел
применяется алгоритм \textbf{GJK} (Gilbert--Johnson--Keerthi).

\paragraph{Минковское разностное множество.}
Для тел $A$ и $B$:
\begin{equation}
  \mathcal{A} \ominus \mathcal{B} =
  \{\vect{a} - \vect{b} \mid \vect{a}\in A,\ \vect{b}\in B\}.
  \label{eq:minkowski}
\end{equation}
Тела пересекаются $\Leftrightarrow$ $\vect{0} \in \mathcal{A}\ominus\mathcal{B}$.

GJK итеративно строит симплекс в $\mathcal{A}\ominus\mathcal{B}$,
приближаясь к началу координат. Если найдено пересечение, вызывается
\textbf{EPA} (Expanding Polytope Algorithm) для глубины проникновения:
\begin{align}
  \delta &= \min_{\vect{x}\in \partial(\mathcal{A}\ominus\mathcal{B})}
  \norm{\vect{x}},
  \label{eq:penetration}\\
  \vect{n} &= \text{нормаль к ближайшей грани политопа EPA}.
  \label{eq:contact_normal}
\end{align}
$\delta$ --- скалярная \textbf{глубина проникновения} (м);
$\vect{n}$ --- единичная нормаль контакта (от $A$ к $B$).

\paragraph{Контакт по треугольнику mesh.}
Для пары «треугольник--выпуклый многогранник» FCL возвращает
1--несколько точек $\vect{c}_k$, нормалей $\vect{n}_k$, глубин $\delta_k$.
При mesh-посадке рамы возможны \textbf{десятки контактов} одновременно
(ребро луча, угол корпуса, несколько треугольников подвеса).

\subsection{Запись контакта в DART}

После narrowphase для пары объектов $(i,j)$ формируется
\textbf{структура контакта} (в Gazebo Classic --- до 250 точек на пару;
в DART --- список произвольной длины, ограничиваемый \code{max\_contacts}).

Для каждой точки контакта $k$:
\begin{equation}
  \mathcal{C}_k =
  \bigl\{
    \text{link}_A,\ \text{link}_B,\;
    \vect{c}_k \in \R^3,\;
    \vect{n}_k \in \R^3,\ \norm{\vect{n}_k}=1,\;
    \delta_k \geq 0,\;
    \vect{v}_{A,k},\ \vect{v}_{B,k}
  \bigr\},
  \label{eq:contact_record}
\end{equation}
\begin{itemize}
  \item $\text{link}_A$, $\text{link}_B$ --- какие звенья SDF участвуют
  (например \code{base\_link} и \code{ground\_plane});
  \item $\vect{c}_k$ --- точка на поверхности в \textbf{мировой} системе;
  \item $\vect{n}_k$ --- нормаль «от $A$ к $B$» (направление запрета проникновения);
  \item $\delta_k$ --- глубина взаимного проникновения (EPA), м;
  \item скорости в точке с учётом вращения звена:
\end{itemize}
\begin{equation}
  \vect{v}_{A,k} = \vect{v}_{A,\mathrm{cm}}
  + \vect{\omega}_A \times (\vect{c}_k - \vect{r}_{A,\mathrm{cm}}),
  \label{eq:contact_velocity}
\end{equation}
аналогично для $B$. Относительная скорость:
$\vect{v}_{\mathrm{rel},k} = \vect{v}_{B,k}-\vect{v}_{A,k}$.

Нормальная и касательные компоненты:
\begin{equation}
  v_{n,k} = \vect{n}_k^\top \vect{v}_{\mathrm{rel},k},
  \qquad
  \vect{v}_{t,k} = \vect{v}_{\mathrm{rel},k} - v_{n,k}\,\vect{n}_k.
  \label{eq:vn}
\end{equation}

\paragraph{Manifold (устойчивость).}
Чтобы контакт не «прыгал» между треугольниками mesh, движки кэшируют
\textbf{манифольд} --- набор $\mathcal{C}_k$, близкий к предыдущему шагу.
Иначе при покое на mesh сила дрожит, хотя визуально дрон неподвижен.

\subsection{Передача энергии при контакте}

Контакт меняет не только силы, но и \textbf{энергию} системы.

\paragraph{Работа контактной силы за шаг.}
\begin{equation}
  W_k^{\mathrm{contact}} \approx
  \vect{F}_{\mathrm{contact},k} \cdot \Delta\vect{p}_k,
  \qquad
  \Delta\vect{p}_k \approx \vect{v}_{\mathrm{rel},k}\,\Delta t,
  \label{eq:contact_work}
\end{equation}
где $\vect{F}_{\mathrm{contact},k}$ из \eqref{eq:contact_force}.
При посадке нормальная составляющая делает отрицательную работу
(тормозит падение), переводя кинетическую энергию в деформацию/импульс.

\paragraph{Импульс и кинетическая энергия.}
Импульс $\vect{J}_k = \vect{F}_{\mathrm{contact},k}\,\Delta t$
изменяет скорости звеньев:
\begin{equation}
  \Delta \vect{v}_A = -\mat{M}_A^{-1}\,\vect{J}_k,\quad
  \Delta \vect{v}_B = +\mat{M}_B^{-1}\,\vect{J}_k
  \quad \text{(упрощённо, без вращения)}.
\end{equation}
Изменение кинетической энергии тела $A$:
\begin{equation}
  \Delta E_{k,A} =
  \tfrac{1}{2} m_A \norm{\vect{v}_A+\Delta\vect{v}_A}^2
  - \tfrac{1}{2} m_A \norm{\vect{v}_A}^2
  \approx \vect{v}_A^\top (m_A \Delta\vect{v}_A).
  \label{eq:delta_KE}
\end{equation}
В полной модели DART учитывается и вращательная часть через
$\vect{J}_k$ и моменты от приложения силы в $\vect{c}_k$.

\paragraph{Нормальный импульс и восстановление.}
Скалярный нормальный импульс $\lambda_{n,k}$ из LCP задаёт скачок
нормальной относительной скорости:
\begin{equation}
  v_{n,k}^{+} = v_{n,k}^{-} + \Delta v_{n,k},
  \qquad
  \Delta v_{n,k} \approx \frac{\lambda_{n,k}}{m_{\mathrm{eff},k}},
  \label{eq:impulse_normal}
\end{equation}
$m_{\mathrm{eff},k}$ --- эффективная масса вдоль $\vect{n}_k$
(из $\vect{J}_k^\top \mat{M}^{-1}\vect{J}_k$).
Коэффициент реституции \eqref{eq:restitution} ограничивает
\textbf{отскок}: часть энергии «теряется» ($e<1$), иначе дрон
бесконечно подпрыгивал бы с пола.

\paragraph{Трение и диссипация.}
Касательный импульс $\vect{\lambda}_{t,k}$ гасит скольжение:
\begin{equation}
  E_{\mathrm{diss},k}^{\mathrm{fric}} \approx
  -\vect{\lambda}_{t,k} \cdot \vect{v}_{t,k},
  \label{eq:fric_diss}
\end{equation}
энергия уходит в тепло (модельно, не считается явно).

\paragraph{Split impulse и «фантомная» энергия.}
Коррекция \eqref{eq:split_impulse} сдвигает \textbf{положение},
не добавляя скорости --- это численный приём против проникновения mesh
в пол, а не физический пружинный удар.

\paragraph{Сравнение с SITL.}
В SITL-only $W_k^{\mathrm{contact}} = 0$: энергия меняется только
от тяги, drag и гравитации \eqref{eq:sitl_v}; удара о пол в уравнениях нет.

\subsection{Контактное ограничение (непроникновение)}

Требование отсутствия проникновения на уровне скоростей:
\begin{equation}
  v_{n,k} + \dot{\delta}_k \geq 0
  \quad\Leftrightarrow\quad
  v_{n,k} \geq -\frac{\delta_k}{\Delta t}
  \label{eq:nonpenetration}
\end{equation}
(в DART используется velocity-based формулировка с допуском Baumgarte).

Импульс $\lambda_{n,k} \geq 0$ вдоль $\vect{n}_k$ удовлетворяет
\textbf{дополнению}:
\begin{equation}
  \lambda_{n,k}\, v_{n,k} = 0.
  \label{eq:complementarity}
\end{equation}
Либо контакт разрывается ($\lambda_{n,k}=0$), либо нормальная
относительная скорость нулевая.

\subsection{Трение Coulomb}

В касательной плоскости с ортонормированным базисом
$\{\vect{t}_{1,k}, \vect{t}_{2,k}\}$:
\begin{equation}
  \vect{\lambda}_{t,k} = \lambda_{t,1,k}\,\vect{t}_{1,k}
  + \lambda_{t,2,k}\,\vect{t}_{2,k},
  \qquad
  \norm{\vect{\lambda}_{t,k}} \leq \mu\,\lambda_{n,k},
  \label{eq:coulomb}
\end{equation}
где $\mu$ --- коэффициент трения из SDF (\code{surface/friction/ode/mu}).

Касательные скорости:
\begin{align}
  v_{t,1,k} &= \vect{t}_{1,k}^\top(\vect{v}_{B,k}-\vect{v}_{A,k}), \\
  v_{t,2,k} &= \vect{t}_{2,k}^\top(\vect{v}_{B,k}-\vect{v}_{A,k}).
\end{align}

\subsection{Матричная форма LCP (boxed)}

Для всех контактов и суставных ограничений собирается система:
\begin{equation}
  \mat{A}\,\vect{\lambda} + \vect{b} = \vect{w},
  \qquad
  \vect{l} \leq \vect{\lambda} \leq \vect{u},
  \qquad
  \vect{\lambda} \circ \vect{w} = \vect{0},
  \label{eq:lcp}
\end{equation}
где $\circ$ --- покомпонентное произведение (complementarity),
$\vect{l},\vect{u}$ --- нижние/верхние границы (для трения --- зависят
от $\lambda_n$ через \emph{findex}).

Блок $\mat{A}$ строится из:
\begin{equation}
  A_{ij} = \vect{J}_i\,\mat{M}^{-1}\,\vect{J}_j^\top,
  \qquad
  b_i = -v_i^{\mathrm{constraint}} + b_i^{\mathrm{bias}},
  \label{eq:lcp_A}
\end{equation}
где $\vect{J}_i$ --- строка якобиана $i$-го ограничения
(нормаль или касательная).

После решения:
\begin{equation}
  \Delta \vect{v} = \mat{M}^{-1}\,\vect{J}^\top\,\vect{\lambda},
  \qquad
  \vect{v}_{n+1} = \vect{v}_n + \Delta\vect{v},
  \label{eq:impulse_apply}
\end{equation}
\begin{equation}
  \vect{F}_{\mathrm{contact},k} \approx
  \frac{\lambda_{n,k}\,\vect{n}_k
  + \lambda_{t,1,k}\,\vect{t}_{1,k}
  + \lambda_{t,2,k}\,\vect{t}_{2,k}}{\Delta t}.
  \label{eq:contact_force}
\end{equation}

\subsection{Коэффициент восстановления (restitution)}

Для упругого отскока вводится целевая нормальная скорость:
\begin{equation}
  v_{n,k}^{\mathrm{target}} = -e\, v_{n,k}^{-},
  \label{eq:restitution}
\end{equation}
$e \in [0,1]$ из \code{surface/bounce/restitution\_coefficient},
$v_{n,k}^{-}$ --- нормальная скорость до импульса.
При $e=0$ --- неупругое касание (мягкая посадка).

\subsection{Split impulse (коррекция позиции)}

Чтобы убрать визуальное проникновение без «взрыва» энергии,
DART/Gazebo применяют \textbf{split impulse}:
\begin{equation}
  \vect{p}_{n+1}^{\mathrm{corr}} =
  \vect{p}_n - \beta\,\sum_k \delta_k\,\vect{n}_k,
  \label{eq:split_impulse}
\end{equation}
где $\beta \in (0,1]$ --- коэффициент коррекции (ERP-подобный параметр).
Скорости при этом корректируются отдельным LCP, что отличается от
простого «подталкивания» тела вверх на $\delta$.

\subsection{Параметры SDF для поверхности}

\begin{center}
\renewcommand{\arraystretch}{1.2}
\begin{tabular}{|l|l|}
\hline
\textbf{SDF-тег} & \textbf{Влияние на модель} \\
\hline
\code{mu}, \code{mu2} & $\mu$ в \eqref{eq:coulomb} \\
\code{slip1}, \code{slip2} & допуск проскальзывания \\
\code{restitution\_coefficient} & $e$ в \eqref{eq:restitution} \\
\code{kp}, \code{kd}, \code{min\_depth} & жёсткость контакта (ODE-совместимые) \\
\code{max\_contacts} & макс.\ число точек на пару тел \\
\hline
\end{tabular}
\end{center}

\subsection{Ограничения суставов (дополнение к LCP)}

Помимо контактов, DART добавляет в \eqref{eq:lcp} строки для:
\begin{itemize}
  \item \textbf{Joint limits:} $q_j \in [q_{j,\min}, q_{j,\max}]$;
  \item \textbf{Joint motors/servo:} $\tau_j = k_p(q_j^{\mathrm{des}}-q_j) - k_d\dot{q}_j$;
  \item \textbf{Weld/fixed:} полное запрещение относительного движения.
\end{itemize}
При нарушении лимита вводится одностороннее ограничение с импульсом,
аналогичным \eqref{eq:complementarity}. Для \code{revolute} ротора
$\theta_i \in [\theta_{\min}, \theta_{\max}]$ (часто без жёсткого лимита,
только скорость через \code{JointVelocityCmd}).

\subsection{Сложность при посадке дрона}

Для квадрокоптера с mesh-коллизией рамы на плоскости $z=0$:
\begin{itemize}
  \item broadphase: $O(m\log m)$;
  \item narrowphase: до $O(|\mathcal{P}_{\mathrm{cand}}| \cdot N_{\triangle})$
  для mesh;
  \item LCP: $O(c^3)$, $c$ --- число контактов + лимитов суставов
  (типично $c \gg 4$ при mesh, против $c=0$ в SITL-only).
\end{itemize}

%====================================================================
\section{ArduPilot SITL: математическая модель}
\label{sec:sitl}
%====================================================================

В проекте \code{Drone\_Swarm\_Simulator\_v2} динамика задаётся классом
\code{SITL::MultiCopter} (\code{SIM\_Multicopter.cpp}) и
\code{Motor::calculate\_forces} (\code{SIM\_Motor.cpp}).

\subsection{Состояние и интегратор}

Вектор состояния (12-мерный):
\[
  \vect{x} = (\vect{p}^\top,\ \vect{\eta}^\top,\ \vect{v}^\top,\ \vect{\omega}^\top)^\top,
\]
NED: $z$ вниз, $g=(0,0,G)^\top$, $G \approx 9.80665$ м/с$^2$.

Уравнения (непрерывный аналог \code{Aircraft::update\_dynamics}):
\begin{align}
  \dot{\vect{p}} &= \vect{v},
  \label{eq:sitl_p}\\
  m\dot{\vect{v}} &= \mat{R}\,\vect{f}_b + m\vect{g},
  \label{eq:sitl_v}\\
  \mat{J}\dot{\vect{\omega}} &= \vect{M}_b - \vect{\omega}\times(\mat{J}\vect{\omega}),
  \label{eq:sitl_omega}\\
  \dot{\mat{R}} &= \mat{R}\,\mat{S}(\vect{\omega}).
  \label{eq:sitl_R}
\end{align}

\textbf{Явное} дискретное обновление (шаг $\Delta t$ из \code{frame\_time\_us}):
\begin{align}
  \vect{\omega} &\leftarrow \vect{\omega} + \dot{\vect{\omega}}\,\Delta t,
  \label{eq:sitl_disc_gyro}\\
  \mat{R} &\leftarrow \mat{R}\,\Delta\mat{R}(\vect{\omega}\Delta t),
  \label{eq:sitl_disc_R}\\
  \vect{v} &\leftarrow \vect{v} + \vect{a}_e\,\Delta t,
  \quad \vect{a}_e = \mat{R}\vect{a}_b + \vect{g},
  \label{eq:sitl_disc_v}\\
  \vect{p} &\leftarrow \vect{p} + \vect{v}\,\Delta t.
  \label{eq:sitl_disc_p}
\end{align}

\subsection{Тяга ротора в SITL (модель импульсного диска)}

Команда из PWM:
\begin{equation}
  u = \mathrm{clip}\Bigl(
  \frac{\mathrm{pwm} - \mathrm{pwm}_{\min}}{\mathrm{pwm}_{\max}-\mathrm{pwm}_{\min}},
  0, 1 \Bigr).
  \label{eq:sitl_pwm}
\end{equation}

Скорость набегания потока на винт (body frame):
\begin{align}
  \vect{v}_{\mathrm{motor}} &= \vect{v}_{\mathrm{air},b}
  - (\vect{r}_{\mathrm{motor}} \times \vect{\omega}),
  \label{eq:sitl_vmotor}\\
  v_{\mathrm{in}} &= \max\bigl(0,\,-\vect{v}_{\mathrm{motor}}\cdot\vect{e}_T\bigr),
  \label{eq:sitl_vin}
\end{align}
где $\vect{e}_T$ --- ось тяги, $\vect{r}_{\mathrm{motor}}$ --- плечо мотора.

Выходная скорость струи и тяга:
\begin{align}
  v_{\mathrm{out}} &= \eta_{\mathrm{volt}}\, v_{\max}\,
  \sqrt{(1-\varepsilon)\,u + \varepsilon\, u^2},
  \label{eq:sitl_vout}\\
  T &= \tfrac{1}{2}\,\rho\, A_{\mathrm{eff}}\,
  \bigl(v_{\mathrm{out}}^2 - v_{\mathrm{in}}^2\bigr),
  \label{eq:sitl_thrust}
\end{align}
где $\rho$ --- плотность воздуха, $A_{\mathrm{eff}}$ --- эффективная площадь
диска, $\varepsilon$ --- \code{mot\_expo}, $\eta_{\mathrm{volt}}$ --- масштаб
от напряжения батареи.

\paragraph{Сравнение с Gazebo.}
В Gazebo \eqref{eq:thrust}: $T \propto \omega^2$ без явного $v_{\mathrm{in}}$.
В SITL \eqref{eq:sitl_thrust} учитывает \textbf{осевой встречный поток}
(inflow), что физичнее для набора высоты, но всё ещё без blade-element theory.

\subsection{Моменты в SITL}

Тяга в body frame: $\vect{F}_T = \vect{e}_T\, T$.
Реактивный момент (yaw), ось вдоль $\vect{e}_T$:
\begin{equation}
  \vect{\tau}_{\mathrm{yaw}} =
  -\zeta_{\mathrm{yaw}}\, u\, T\,\vect{e}_T,
  \label{eq:sitl_yaw}
\end{equation}
$\zeta_{\mathrm{yaw}} = 0.05 \cdot d_{\mathrm{diag}} \cdot \code{yaw\_factor}$
(из \code{SIM\_Motor.cpp}, $d_{\mathrm{diag}}$ --- размер рамы).

Полный момент:
\begin{equation}
  \vect{M}_b = \vect{r}_{\mathrm{motor}} \times \vect{F}_T
  + \vect{\tau}_{\mathrm{yaw}} + \cdots
  \label{eq:sitl_moment}
\end{equation}

\subsection{Momentum drag (опционально)}

При \code{use\_drag=true}:
\begin{equation}
  \vect{F}_{\mathrm{md}} =
  -c_{\mathrm{md}}\sqrt{\rho}\,
  \begin{pmatrix}
    v_x (\sqrt{|F_{T,y}|}+\sqrt{|F_{T,z}|}) \\
    v_y (\sqrt{|F_{T,x}|}+\sqrt{|F_{T,z}|}) \\
    v_z (\sqrt{|F_{T,x}|}+\sqrt{|F_{T,y}|}+\sqrt{|F_{T,z}|})
  \end{pmatrix},
  \label{eq:sitl_momentum_drag}
\end{equation}
$c_{\mathrm{md}} = \code{momentum\_drag\_coefficient}$.
В Gazebo аналог --- \eqref{eq:blade_drag}, но с другой структурой.

\subsection{Что SITL \emph{не} считает}

\begin{itemize}
  \item контакты с землёй, стенами, другими дронами;
  \item отдельные вращающиеся звенья роторов в дереве кинематики;
  \item LCP и $\mat{M}(\vect{q})$ полной системы;
  \item mesh-коллизии и множественные точки контакта;
  \item joint limits, friction в осях (кроме параметров рамы).
\end{itemize}
Высота в SITL-only ограничена только управлением автопилота, не физикой среды.

%====================================================================
\section{Сравнение SITL и Gazebo: сложность модели}
\label{sec:compare}
%====================================================================

\subsection{Сводная таблица}

\begin{center}
\small
\renewcommand{\arraystretch}{1.25}
\begin{tabular}{|p{3.8cm}|p{5.2cm}|p{5.5cm}|}
\hline
\textbf{Аспект} & \textbf{ArduPilot SITL-only} & \textbf{Gazebo Sim + DART} \\
\hline
Число тел & 1 жёсткое тело (рамка) & Дерево: корпус + 4 ротора + среда \\
\hline
Степени свободы & 6 (фикс.\ тело) & $6 + n_{\mathrm{joint}}$ (тип.\ $6+4$) \\
\hline
Уравнения & \eqref{eq:sitl_p}--\eqref{eq:sitl_R} & \eqref{eq:dart_eom} + \eqref{eq:lcp} \\
\hline
Матрица масс & постоянная $\mat{J}$, $m$ & $\mat{M}(\vect{q})$ полная \\
\hline
Интегратор & явный Эйлер \eqref{eq:sitl_disc_v} & semi-implicit + LCP \\
\hline
Тяга & \eqref{eq:sitl_thrust} (inflow) & \eqref{eq:thrust} ($\omega^2$) \\
\hline
Контакты & \textbf{нет} & \eqref{eq:contact_force}, \eqref{eq:coulomb} \\
\hline
Mesh / сетка & не используется & GJK/EPA, convex decomp. \\
\hline
Параметры среды & $\rho$, $g$, батарея & $\rho$, $g$, $\mu$, $e$, mesh, wind \\
\hline
Стоимость шага & $O(1)$ по телам & $O(n) + O(c^3)$ контакты \\
\hline
\end{tabular}
\end{center}

Ниже --- разбор \textbf{каждой строки} с явными формулами.

%--------------------------------------------------------------------
\subsection{Число тел}
%--------------------------------------------------------------------

\textbf{SITL:} одно агрегированное тело с массой $m$ и инерцией $\mat{J}$.
Все четыре мотора --- точки приложения сил:
\begin{equation}
  \vect{F}_b = \sum_{i=1}^{4} \vect{F}_{T,i}, \qquad
  \vect{M}_b = \sum_{i=1}^{4}
  \bigl(\vect{r}_i \times \vect{F}_{T,i} + \vect{\tau}_{Q,i}\bigr).
  \label{eq:sitl_aggregate}
\end{equation}
Геометрических «тел» для контакта нет.

\textbf{Gazebo:} минимум $n_{\mathrm{body}} = 5$ звеньев на дрон
(1 корпус + 4 ротора) плюс $n_{\mathrm{env}}$ тел среды (пол, препятствия):
\begin{equation}
  n_{\mathrm{total}} = n_{\mathrm{body}} + n_{\mathrm{env}},
  \qquad
  n_{\mathrm{coll}} \geq n_{\mathrm{total}}
  \quad \text{(если у link несколько \code{collision})}.
  \label{eq:nbodies}
\end{equation}
Каждое звено --- отдельный \code{BodyNode} в \eqref{eq:dart_eom};
контакт «корпус--пол» и «луч--стена» --- разные пары в
\eqref{eq:broadphase}.

%--------------------------------------------------------------------
\subsection{Степени свободы}
%--------------------------------------------------------------------

\textbf{SITL:} вектор состояния размерности 12:
\begin{equation}
  \dim \vect{x}_{\mathrm{SITL}} = 12,
  \quad
  \vect{x} = (p_x,p_y,p_z,\eta_x,\eta_y,\eta_z,
  v_x,v_y,v_z,\omega_x,\omega_y,\omega_z)^\top.
  \label{eq:sitl_state_dim}
\end{equation}
Три поступательных + три вращательных DOF \emph{всего аппарата};
скорости роторов не входят в $\vect{x}$.

\textbf{Gazebo:} при \code{free} + 4$\times$\code{revolute}:
\begin{equation}
  \dim \vect{q} = 6 + 4 = 10,
  \qquad
  \dim \dot{\vect{q}} = 10,
  \label{eq:gazebo_dof}
\end{equation}
плюс скрытые ограничения контактов (не координаты, а $\vect{\lambda}$
в \eqref{eq:lcp}). При фиксированном корпусе и \code{revolute} только
на моторах: $\dim\vect{q} = 4$, но положение корпуса всё равно 6 DOF
в мире через \code{free joint}.

%--------------------------------------------------------------------
\subsection{Уравнения движения}
%--------------------------------------------------------------------

\textbf{SITL} (замкнутая система без контактов):
\begin{equation}
  \dot{\vect{x}} =
  f_{\mathrm{SITL}}(\vect{x}, \vect{u}_{\mathrm{PWM}}, \rho, V_{\mathrm{bat}}),
  \label{eq:sitl_ode}
\end{equation}
где $f_{\mathrm{SITL}}$ реализует \eqref{eq:sitl_p}--\eqref{eq:sitl_R}
с силами \eqref{eq:sitl_thrust}--\eqref{eq:sitl_moment}.

\textbf{Gazebo:}
\begin{equation}
  \underbrace{\mat{M}(\vect{q})\ddot{\vect{q}}
  + \vect{C} + \vect{G}}_{\text{динамика скелета}}
  =
  \underbrace{\vect{\tau}_{\mathrm{motor}}}_{\text{плагины}}
  +
  \underbrace{\vect{J}_c^\top \vect{\lambda}}_{\text{контакты}},
  \label{eq:gazebo_full}
\end{equation}
плюс алгебраические ограничения \eqref{eq:lcp}.
Размерность \eqref{eq:gazebo_full}: $n_q$ уравнений + $c$ неравенств LCP.

%--------------------------------------------------------------------
\subsection{Матрица масс}
%--------------------------------------------------------------------

\textbf{SITL:} постоянные скаляр $m$ и тензор $\mat{J}$ (из параметров рамы):
\begin{equation}
  m\dot{\vect{v}} = \mat{R}\vect{f}_b + m\vect{g},
  \qquad
  \mat{J}\dot{\vect{\omega}} = \vect{M}_b - \vect{\omega}\times\mat{J}\vect{\omega}.
  \label{eq:sitl_mass_const}
\end{equation}
$\mat{M}_{\mathrm{SITL}} = \mathrm{diag}(m\mat{E}_3, \mat{J})$ в 6D-пространстве
(в связанных осях для вращения).

\textbf{Gazebo:} полная $\mat{M}(\vect{q}) \in \R^{n_q \times n_q}$:
\begin{equation}
  M_{ij}(\vect{q}) =
  \frac{\partial}{\partial \ddot{q}_j}
  \Bigl[\tau_i^{\mathrm{ID}}(\vect{q},\dot{\vect{q}},\ddot{\vect{q}})\Bigr],
  \label{eq:mass_matrix_def}
\end{equation}
учитывает:
\begin{itemize}
  \item массу корпуса и каждого ротора;
  \item положение $\theta_i$ (меняется $\mat{M}$ при повороте лучей);
  \item связи суставов (недиагональные $M_{ij}$).
\end{itemize}
Пример: при разгоне $\dot{\theta}_i$ ротор тянет корпус --- в SITL это
только через $\vect{F}_b$; в DART ещё через $M_{ij}$ между $\ddot{\theta}_i$
и $\ddot{p}$.

%--------------------------------------------------------------------
\subsection{Интегратор}
%--------------------------------------------------------------------

\textbf{SITL --- явный (forward) Euler.}
Из \eqref{eq:sitl_disc_gyro}--\eqref{eq:sitl_disc_p}:
\begin{align}
  \vect{\omega}_{n+1} &= \vect{\omega}_n + \ddot{\vect{\omega}}_n\,\Delta t,
  \\
  \mat{R}_{n+1} &= \mat{R}_n\,\Delta\mat{R}(\vect{\omega}_n\Delta t),
  \\
  \vect{v}_{n+1} &= \vect{v}_n + (\mat{R}_n\vect{a}_{b,n} + \vect{g})\,\Delta t,
  \\
  \vect{p}_{n+1} &= \vect{p}_n + \vect{v}_n\,\Delta t.
  \label{eq:sitl_euler_explicit}
\end{align}
Замечание: скорость в \eqref{eq:sitl_disc_p} берётся \textbf{старой}
$\vect{v}_n$ (не $\vect{v}_{n+1}$) --- это чистый явный Эйлер;
при больших $\Delta t$ возможна нестабильность (в SITL $\Delta t$ малый,
$\sim 1/300$ с).

\textbf{Gazebo --- semi-implicit Euler + LCP.}
Сначала решают импульсы при \emph{текущих} скоростях, затем:
\begin{align}
  \dot{\vect{q}}_{n+1} &= \dot{\vect{q}}_n + \ddot{\vect{q}}_n\,\Delta t,
  \label{eq:gazebo_vel_int}\\
  \vect{q}_{n+1} &= \vect{q}_n + \dot{\vect{q}}_{n+1}\,\Delta t,
  \label{eq:gazebo_pos_int}
\end{align}
где $\ddot{\vect{q}}_n$ уже включает $\vect{J}_c^\top\vect{\lambda}$ из
\eqref{eq:impulse_apply}. Это \textbf{симплектическая} схема для
свободного движения; контакт добавляет диссипацию \eqref{eq:fric_diss}.

\paragraph{Сравнение устойчивости.}
При контакте явный SITL без пола не решает LCP; Gazebo на каждом шаге
при $\delta_k > 0$ решает \eqref{eq:lcp} --- иначе тело провалится сквозь
\code{plane}.

%--------------------------------------------------------------------
\subsection{Тяга}
%--------------------------------------------------------------------

\textbf{SITL:} \eqref{eq:sitl_thrust} с inflow \eqref{eq:sitl_vin}.

\textbf{Gazebo:} \eqref{eq:thrust}--\eqref{eq:blade_drag} с фильтром
\eqref{eq:motor_filter}. Эквивалентность при $v_{\mathrm{in}} \approx 0$,
$\rho=\mathrm{const}$: $T \approx k_f \omega^2$ с подбором $k_f$.

%--------------------------------------------------------------------
\subsection{Контакты}
%--------------------------------------------------------------------

\textbf{SITL:}
\begin{equation}
  \vect{F}_{\mathrm{contact}}^{\mathrm{SITL}} \equiv \vect{0},
  \qquad
  \vect{\lambda} \equiv \vect{0}.
  \label{eq:sitl_no_contact}
\end{equation}

\textbf{Gazebo:} при $c$ контактах полная система \eqref{eq:lcp}--\eqref{eq:contact_force};
энергия \eqref{eq:delta_KE}--\eqref{eq:fric_diss}.

%--------------------------------------------------------------------
\subsection{Mesh и детектор}
%--------------------------------------------------------------------

\textbf{SITL:} нет \eqref{eq:broadphase}--\eqref{eq:penetration}.

\textbf{Gazebo:} для каждой пары из $\mathcal{P}_{\mathrm{cand}}$
GJK $\Rightarrow$ при пересечении EPA $\Rightarrow$ $\delta_k$, $\vect{n}_k$
$\Rightarrow$ $\mathcal{C}_k$ \eqref{eq:contact_record}.

%--------------------------------------------------------------------
\subsection{Параметры среды}
%--------------------------------------------------------------------

\textbf{SITL:} $\rho$ в \eqref{eq:sitl_thrust}, $\vect{g}$ в \eqref{eq:sitl_v},
$V_{\mathrm{bat}}$ в \eqref{eq:sitl_vout}; $\mu$, $e$, mesh --- не определены.

\textbf{Gazebo:} дополнительно $\mu$ \eqref{eq:coulomb}, $e$ \eqref{eq:restitution},
геометрия пола, \code{Wind} в \eqref{eq:vrel}.

%--------------------------------------------------------------------
\subsection{Стоимость шага}
%--------------------------------------------------------------------

\textbf{SITL:}
\begin{equation}
  T_{\mathrm{step}}^{\mathrm{SITL}} =
  O(1)\ \text{(4 мотора)} + O(1)\ \text{(интеграция 12 состояний)}.
  \label{eq:cost_sitl}
\end{equation}

\textbf{Gazebo:}
\begin{equation}
  T_{\mathrm{step}}^{\mathrm{GZ}} =
  O(n)\ \text{(ABA)} +
  O(m \log m)\ \text{(broadphase)} +
  |\mathcal{P}_{\mathrm{cand}}|\,T_{\mathrm{GJK/EPA}} +
  O(c^3)\ \text{(LCP)}.
  \label{eq:cost_gazebo}
\end{equation}
При посадке mesh-рамы $c$ может быть 10--100+, доминирует LCP.

\subsection{Почему Gazebo «сложнее» при тех же выходных величинах}

С точки зрения пользователя автопилота важны
$\vect{p}, \mat{R}, \vect{v}, \vect{\omega}$.
И SITL, и Gazebo их выдают, но:

\begin{enumerate}
  \item \textbf{Размерность задачи.}
  SITL: 12 ОДУ + алгебра тяг. Gazebo: столько же \emph{плюс}
  обобщённые координаты суставов, $\mat{M}(\vect{q})$, $\vect{C}$, $\vect{G}$,
  и при посадке --- сотни неравенств \eqref{eq:lcp}.

  \item \textbf{Столкновения.}
  В Gazebo посадка --- решение контактной задачи с трением
  \eqref{eq:coulomb}--\eqref{eq:impulse_apply}.
  В SITL-only дрон может «провалиться» сквозь пол: физического пола нет.

  \item \textbf{Геометрия.}
  Gazebo: mesh $\Rightarrow$ треугольники $\Rightarrow$ множество
  $\vect{c}_k, \delta_k$ в \eqref{eq:contact_record}.
  SITL: геометрический объём аппарата не участвует в динамике.

  \item \textbf{Связность.}
  DART учитывает вращение пропеллеров как DOF суставов
  (\eqref{eq:joint_vel}); SITL --- кинематически не моделирует отдельные
  лопасти, только агрегированную тягу.

  \item \textbf{Численная устойчивость.}
  Gazebo вынужден решать LCP каждый шаг при контакте;
  SITL --- простое суммирование сил без complementarity.
\end{enumerate}

\subsection{Схема потоков данных}

\begin{center}
\begin{tabular}{|c|c|}
\hline
\textbf{SITL-only (Kursov3)} & \textbf{Gazebo Sim} \\
\hline
PWM $\to$ \eqref{eq:sitl_pwm} $\to$ \eqref{eq:sitl_thrust}
& Actuators $\to$ \eqref{eq:motor_filter} $\to$ \eqref{eq:thrust} \\
$\to$ \eqref{eq:sitl_moment} $\to$ \eqref{eq:sitl_v}--\eqref{eq:sitl_omega}
& $\to$ \eqref{eq:newton}, \eqref{eq:euler} $\to$ \eqref{eq:dart_eom} \\
$\to$ \eqref{eq:sitl_disc_p}
& $\to$ collision $\to$ \eqref{eq:lcp} $\to$ \eqref{eq:dart_pos} \\
\hline
\end{tabular}
\end{center}

\subsection{Режим Webots в проекте}

При \code{SIM\_Webots\_Python} интегрирование выполняет \textbf{Webots}
(свой физический движок и контакты); SITL только принимает
\code{fdm\_packet}. Это третий вариант: ближе к Gazebo по наличию среды,
но формулы зависят от движка Webots, а не DART.

%====================================================================
\section{Численный цикл одного шага симуляции}
%====================================================================

Упрощённая последовательность (\code{Physics::Update} + systems
\code{PreUpdate}):

\begin{enumerate}
  \item Системы \code{PreUpdate} (в т.\,ч.\ \code{MulticopterMotorModel}):
  читают команды, выставляют \code{JointVelocityCmd}, добавляют
  \code{ExternalWorldWrenchCmd} (силы/моменты).
  \item \code{Physics}: DART \code{World::step()}:
  \begin{enumerate}
    \item detect collisions;
    \item собрать joint + contact constraints;
    \item решить LCP $\Rightarrow$ импульсы;
    \item интегрировать $\vect{q}, \vect{v}, \vect{\omega}$ за $\Delta t$.
  \end{enumerate}
  \item Системы \code{PostUpdate}: публикация одometry, sensors.
\end{enumerate}

Параметры мира из \code{quadcopter.sdf}:
\code{max\_step\_size} = 0.001 с, \code{real\_time\_factor} = 1.

%====================================================================
\section{Пример численных параметров (модель X3)}
%====================================================================

Из \code{examples/worlds/quadcopter.sdf} для каждого ротора:
\begin{align*}
  k_f &= 8.54858\times 10^{-6}\ \text{Н/(рад/с)}^2, &
  k_m &= 0.016\ \text{м}, \\
  \Omega_{\max} &= 800\ \text{рад/с}, &
  n_s &= 10, \\
  \tau_\uparrow &= 0.0125\ \text{с}, &
  \tau_\downarrow &= 0.025\ \text{с}, \\
  b_f &= 8.06428\times 10^{-5}, &
  \mu_1 &= 10^{-6}.
\end{align*}
Максимальная тяга одного ротора (из \eqref{eq:thrust}):
\[
  T_{\max} \approx k_f\,\Omega_{\max}^2
  \approx 8.55\times 10^{-6} \times 800^2 \approx 5.47\ \text{Н}.
\]
Для 4 роторов $T_{\Sigma,\max} \approx 21.9$ Н --- порядок массы X3
($m \approx 1.5$--2 кг) при $g \approx 9.81$ м/с$^2$.

%====================================================================
\section{Ссылки на исходный код}
%====================================================================

\begin{itemize}
  \item Gazebo Sim: \url{https://github.com/gazebosim/gz-sim}
  \item DART: \url{https://github.com/dartsim/dart},
  \code{docs/onboarding/architecture.md}
  \item FCL (GJK/EPA): \url{https://github.com/flexible-collision-library/fcl}
  \item ArduPilot SITL: \code{libraries/SITL/SIM\_Motor.cpp},
  \code{SIM\_Aircraft.cpp}, \code{SIM\_Multicopter.cpp}
  \item Gazebo: \code{MulticopterMotorModel.cc}, \code{LiftDrag.cc},
  \code{Physics.cc}, \code{examples/worlds/quadcopter.sdf}
\end{itemize}

%====================================================================
\section*{Заключение}
%====================================================================

\textbf{ArduPilot SITL} в режиме проекта Kursov3 решает компактную задачу:
12 уравнений \eqref{eq:sitl_p}--\eqref{eq:sitl_R} с тягой
\eqref{eq:sitl_thrust} и без контактов.
\textbf{Gazebo Sim + DART} добавляет:
\begin{enumerate}
  \item полную динамику связной системы \eqref{eq:dart_eom};
  \item неявную интеграцию с LCP \eqref{eq:lcp}--\eqref{eq:impulse_apply};
  \item обнаружение столкновений \eqref{eq:broadphase}--\eqref{eq:penetration}
  для примитивов и mesh-сеток;
  \item контактные силы, трение \eqref{eq:coulomb} и восстановление
  \eqref{eq:restitution}.
\end{enumerate}

Для дипломного сравнения: при одинаковых командах моторов
\textit{кинематика в воздухе} может быть близкой (обе модели --- 6-DOF +
квадратичная/импульсная тяга), но \textit{взаимодействие со средой}
(посадка, удар, трение, несколько дронов) в Gazebo описывается
существенно более богатой математикой раздела~\ref{sec:collision}
и~\ref{sec:compare}, тогда как SITL-only эти эффекты не моделирует.

\end{document}
